% !PRINTED 767
\weberpage{767}
Durch 5. und 6. ist aber nachgewiesen, dass, wenn $p$ gross genug ist, die Gleichung (15) nicht bestehen kann. Dadurch ist 4. und damit das ganze Theorem I. bewiesen.

Dieser Satz gestattet nun mannigfaltige Anwendungen. Er giebt uns zunächst die Transcendenz von $e$, wenn wir $C_1, C_2, \ldots, z_1, z_2, \ldots$ als ganze rationale Zahlen annehmen. Er giebt ferner die Transcendenz von $\pi$. Denn aus der Gleichung $1+e^{i\pi}=0$ folgt nach I., dass $i\pi$ und folglich auch $\pi$ nicht algebraisch sein kann.

Es folgt ferner daraus:

\gesperrt{Für jede algebraische Zahl $x$, mit Ausnahme von $x=0$, ist $X=e^x$ eine transcendente Zahl.}

\gesperrt{Für jedes algebraische $X$, mit Ausnahme von $X=1$, ist jeder natürliche Logarithmus $x=\log X$ eine transcendente Zahl.}

\gesperrt{Für jeden Bogen, der zum Radius in einem algebraisch ausdrückbaren Verhältnisse $x$ steht, mit Ausnahme von $x=0$, ist $X=\sin x$ eine transcendente Zahl.}

Dies folgt nach I. aus $2i\,X=e^{ix}-e^{-ix}$.

% [sic] -- the source prints the chord as $\frac12\sin(x/2)$; for unit radius the chord is $2\sin(x/2)$.
% [sic] -- the source's claim for $\mathrm{tg}\,x=\alpha x$ omits the algebraic root $x=0$; the intended claim concerns nonzero roots.
Dasselbe gilt für die anderen trigonometrischen Linien, $\cos x$, $\mathrm{tg}\,x$ und für die Sehne $\dfrac{1}{2}\sin\dfrac{x}{2}$. Um noch eins anzuführen: Die transcendente Gleichung $\mathrm{tg}\,x=\alpha x$ hat für ein algebraisches $\alpha$ ausser 0 nur transcendente Zahlen zu Wurzeln.

\begin{center}
\rule{4cm}{0.4pt}
\end{center}
